%======================================================================
%  Capitulo 4 - Diseno logico de bases de datos
%======================================================================

\section{Diseño lógico de bases de datos}

El diseño lógico transforma el modelo conceptual ajustado del módulo de
operación de envíos (figura~\ref{fig:dmlogico}; en notación de Chen,
figura~\ref{fig:erfinal}) en un modelo relacional:
tablas, columnas, claves primarias y claves foráneas, con sus reglas de
integridad. El resultado tiene 19 tablas, 130 columnas y 27 claves
foráneas.

%-----------------------------------------------------------------------
\subsection{Conversión del diseño conceptual al diseño lógico}

\subsubsection{Modelo Entidad Relación a Modelo Relacional}

La conversión aplica las tres reglas básicas del curso y dos reglas
complementarias para los casos que el modelo del módulo incluye:

\begin{enumerate}[label=\textbf{R\arabic*.},leftmargin=1.2cm,itemsep=2pt]
  \item \textbf{Toda entidad se convierte en tabla}, y su identificador
        en clave primaria. Las 19 entidades dan 19 tablas.
  \item \textbf{Toda relación 1:N se convierte en una propagación de
        clave}: la clave primaria del lado ``uno'' pasa como clave
        foránea a la tabla del lado ``muchos''. Si el mínimo del lado
        ``uno'' es 0, la clave foránea admite nulos.
  \item \textbf{Toda relación N:M se convierte en una tabla intermedia.}
        El módulo no tiene relaciones N:M: la única del prediseño
        (\textit{se asigna}) ya se convirtió en la entidad
        ASIGNACION\_DESPACHO en el ajuste 2.
  \item \textbf{Una relación 1:1 se convierte en una clave foránea única}
        en la tabla de participación opcional (PRUEBA\_ENTREGA y
        NOVEDAD respecto a ASIGNACION\_DESPACHO).
  \item \textbf{Una especialización se convierte en una tabla por
        subtipo}, cuya clave primaria es también clave foránea hacia la
        tabla del supertipo (MENSAJERO respecto a EMPLEADO).
\end{enumerate}

La tabla~\ref{tab:conversion} muestra cómo se convirtió cada relación.

{\singlespacing\small
\begin{longtable}{L{2.0cm} L{1.0cm} L{4.9cm} L{4.6cm} L{1.0cm}}
\caption{Conversión de las relaciones del modelo ER del módulo.}
\label{tab:conversion}\\
\toprule
\textbf{Relación} & \textbf{Regla} & \textbf{Entidades (lado uno -- lado muchos)} & \textbf{Clave foránea resultante} & \textbf{Nulo} \\
\midrule
\endfirsthead
\toprule
\textbf{Relación} & \textbf{Regla} & \textbf{Entidades (lado uno -- lado muchos)} & \textbf{Clave foránea resultante} & \textbf{Nulo} \\
\midrule
\endhead
\bottomrule
\endfoot
\textit{cubre} & R2 & CENTRO\_DISTRIBUCION -- ZONA & zona.id\_centro & No \\
\textit{labora} & R2 & CENTRO\_DISTRIBUCION -- EMPLEADO & empleado.id\_centro & No \\
\textit{dispone} & R2 & CENTRO\_DISTRIBUCION -- VEHICULO & vehiculo.id\_centro & No \\
\textit{origina} & R2 & CENTRO\_DISTRIBUCION -- DESPACHO & despacho.id\_centro & No \\
\textit{ubica} & R2 & ZONA -- SEDE\_CLIENTE & sede\_cliente.id\_zona & No \\
\textit{destina} & R2 & ZONA -- ENVIO & envio.id\_zona & No \\
\textit{atiende} & R2 & ZONA -- MENSAJERO & mensajero.id\_zona & Sí \\
\textit{tiene} & R2 & CLIENTE -- SEDE\_CLIENTE & sede\_cliente.id\_cliente & No \\
\textit{radica} & R2 & CLIENTE -- RECLAMACION & reclamacion.id\_cliente & No \\
\textit{clasifica} & R2 & MODALIDAD -- ENVIO & envio.id\_modalidad & No \\
\textit{solicita} & R2 & SEDE\_CLIENTE -- SOLICITUD\_RECOLECCION & solicitud\_recoleccion.id\_sede & No \\
\textit{contiene} & R2 & SOLICITUD\_RECOLECCION -- ENVIO & envio.id\_solicitud & No \\
\textit{recoge} & R2 & DESPACHO -- SOLICITUD\_RECOLECCION & solicitud\_recoleccion.id\_despacho & Sí \\
\textit{dirigido} & R2 & DESTINATARIO -- ENVIO & envio.id\_destinatario & No \\
\textit{se asigna} & R2 & ENVIO -- ASIGNACION\_DESPACHO & asignacion\_despacho.numero\_guia & No \\
\textit{agrupa} & R2 & DESPACHO -- ASIGNACION\_DESPACHO & asignacion\_despacho.id\_despacho & No \\
\textit{soporta} & R4 & ASIGNACION\_DESPACHO -- PRUEBA\_ENTREGA & prueba\_entrega.id\_asignacion (única) & No \\
\textit{presenta} & R4 & ASIGNACION\_DESPACHO -- NOVEDAD & novedad.id\_asignacion (única) & No \\
\textit{registra} & R2 & ENVIO -- EVENTO\_ENVIO & evento\_envio.numero\_guia & No \\
\textit{indica} & R2 & ESTADO\_ENVIO -- EVENTO\_ENVIO & evento\_envio.id\_estado & No \\
\textit{reporta} & R2 & EMPLEADO -- EVENTO\_ENVIO & evento\_envio.id\_empleado & No \\
\textit{refiere} & R2 & ENVIO -- RECLAMACION & reclamacion.numero\_guia & Sí \\
\textit{usa} & R2 & VEHICULO -- DESPACHO & despacho.id\_vehiculo & No \\
\textit{realiza} & R2 & MENSAJERO -- DESPACHO & despacho.id\_mensajero & No \\
\textit{categoriza} & R2 & TIPO\_VEHICULO -- VEHICULO & vehiculo.id\_tipo\_vehiculo & No \\
\textit{conduce} & R2 & TIPO\_VEHICULO -- MENSAJERO & mensajero.id\_tipo\_vehiculo & No \\
\textit{es un} & R5 & EMPLEADO -- MENSAJERO & mensajero.id\_empleado (PK y FK) & No \\
\end{longtable}
\par}

Las claves foráneas que admiten nulos son exactamente las de las
relaciones con mínimo 0 en el lado ``uno'' (\textit{atiende},
\textit{recoge} y \textit{refiere}); las demás son obligatorias.

\subsubsection{Reglas de Integridad (Propias del Modelo relacional y propias del desarrollador)}

\textbf{Reglas propias del modelo relacional.} Se garantizan con la
estructura de las tablas (tabla~\ref{tab:integridad}).

{\singlespacing\small
\begin{longtable}{L{3.2cm} L{8.2cm} L{2.6cm}}
\caption{Reglas de integridad propias del modelo relacional.}
\label{tab:integridad}\\
\toprule
\textbf{Regla} & \textbf{Cómo se aplica en el módulo} & \textbf{Restricciones} \\
\midrule
\endfirsthead
\toprule
\textbf{Regla} & \textbf{Cómo se aplica en el módulo} & \textbf{Restricciones} \\
\midrule
\endhead
\bottomrule
\endfoot
Integridad de entidad & Cada tabla tiene clave primaria, única y no nula. Se usan claves sustitutas numéricas, salvo ENVIO, cuya clave es el número de guía. Las claves naturales (NIT, documento, placa, nombres de catálogos) se declaran únicas. & 19 PRIMARY KEY, 13 UNIQUE \\
\midrule
Integridad referencial & Toda clave foránea apunta a una fila existente. No se define borrado en cascada: la política de acceso prohíbe eliminar registros, que se anulan por estado, así que el motor rechaza borrar una fila referenciada. & 27 FOREIGN KEY \\
\midrule
Integridad de dominio & Cada columna tiene tipo y tamaño; las obligatorias son NOT NULL; los estados y tipos solo admiten los valores de su lista, y los estados iniciales tienen valor por defecto. & 110 NOT NULL, 33 CHECK, 12 DEFAULT \\
\end{longtable}
\par}

\textbf{Reglas propias del desarrollador.} Son las reglas de negocio
que salen de las políticas y de la caracterización de los procesos
(tabla~\ref{tab:reglasnegocio}). Las que se pueden expresar sobre una
sola fila o con unicidad ya están en el modelo lógico; las que
dependen de otras filas o tablas se implementarán con disparadores,
procedimientos y permisos en el diseño físico.

{\singlespacing\small
\begin{longtable}{L{1.0cm} L{6.6cm} L{2.6cm} L{4.2cm}}
\caption{Reglas de negocio del módulo.}
\label{tab:reglasnegocio}\\
\toprule
\textbf{Cód.} & \textbf{Regla} & \textbf{Origen} & \textbf{Implementación} \\
\midrule
\endfirsthead
\toprule
\textbf{Cód.} & \textbf{Regla} & \textbf{Origen} & \textbf{Implementación} \\
\midrule
\endhead
\bottomrule
\endfoot
RN01 & Cada envío tiene un número de guía único. & Calidad del servicio & PK de envio \\
RN02 & Todo envío se registra con destinatario, dirección, zona, modalidad, peso, contenido y fecha de compromiso. & Calidad del servicio & NOT NULL en envio \\
RN03 & La fecha de compromiso no es anterior a la de registro, los pesos y el flete son positivos y el valor declarado no es negativo. & P1 & CHECK envio\_compromiso\_ck, envio\_pesos\_ck, envio\_valor\_ck \\
RN04 & Un destinatario existe una sola vez (tipo y número de documento). & Datos personales & UNIQUE destinatario\_documento\_uk \\
RN05 & Ningún envío se marca como entregado sin prueba de entrega. & Calidad del servicio & Disparador (diseño físico) \\
RN06 & Los envíos frágiles o de alto valor solo se entregan al destinatario: el documento de quien recibe debe ser el del destinatario. & Manejo de la mercancía & Disparador (diseño físico) \\
RN07 & Máximo dos intentos de entrega; el tercero solo puede ser la devolución y solo después de dos novedades. & P3 & CHECK asignacion\_intento\_ck, UNIQUE asignacion\_intento\_uk y disparador \\
RN08 & Cada intento tiene como máximo una prueba de entrega o una novedad. & P2 & UNIQUE prueba\_asignacion\_uk, novedad\_asignacion\_uk \\
RN09 & Solo un mensajero puede realizar un despacho. & Seguridad vial & FK despacho\_mensajero\_fk \\
RN10 & Al crear un despacho, el vehículo debe estar DISPONIBLE y el mensajero ACTIVO (la validación de documentos, licencia y capacitación se agrega con M3 y M4). & Seguridad vial & Disparador (diseño físico) \\
RN11 & El vehículo del despacho debe ser del tipo que conduce el mensajero. & Seguridad vial & Disparador (diseño físico) \\
RN12 & Todo envío nace con el evento REGISTRADO y sus estados siguen el ciclo de vida (no se pasa a ENTREGADO sin haber estado EN\_RUTA). & Seguimiento por eventos & Procedimiento de registro y disparador \\
RN13 & Cada evento registra el empleado responsable. & Manejo de la mercancía & FK NOT NULL evento\_empleado\_fk \\
RN14 & La respuesta a una reclamación es posterior a su radicación; la meta de cuatro horas se mide como indicador. & OE2 & CHECK reclamacion\_resp\_ck y vista \\
RN15 & Solo se registran solicitudes de clientes activos (hasta integrar los contratos de M2). & P1 & Disparador (diseño físico) \\
RN16 & No se eliminan registros; se anulan o cierran por estado. & Acceso a la información & Permisos sin DELETE (diseño físico) \\
\end{longtable}
\par}

\subsubsection{Optimización del Diseño Lógico de la BD}

Antes de normalizar se tomaron las siguientes decisiones de diseño:

\begin{enumerate}[label=\arabic*.,leftmargin=1.2cm,itemsep=2pt]
  \item \textbf{Claves.} El número de guía es la clave de ENVIO porque
        es estable, va impreso en el paquete y es lo que el cliente
        consulta. Las demás tablas usan claves sustitutas numéricas, más
        cortas y estables que las naturales, que se conservan como
        únicas.
  \item \textbf{Estado actual del envío.} No se guarda en ENVIO, porque
        se deduce del último evento y guardarlo dos veces permitiría
        contradicciones (ajuste 4). Se obtiene con la vista
        v\_estado\_actual\_envio; si en el diseño físico se necesita por
        rendimiento, se agregará como dato derivado mantenido por un
        disparador.
  \item \textbf{Prueba y novedad en tablas propias.} Podrían ser columnas
        de ASIGNACION\_DESPACHO, pero cada intento tiene una u otra, nunca
        las dos, y la mayoría de columnas quedaría nula. Separarlas con
        una clave foránea única evita los nulos y permite reglas propias
        para cada una.
  \item \textbf{Mensajero como tabla de subtipo.} Guardar la zona y el
        tipo de vehículo en EMPLEADO dejaría esas columnas vacías para
        quien no es mensajero, y no impediría asignar un despacho a otro
        cargo. La tabla MENSAJERO resuelve ambas cosas.
  \item \textbf{Catálogos frente a listas fijas.} Son tablas los
        dominios que tienen atributos propios o que usan las reglas
        (estados del envío, modalidades, tipos de vehículo); son listas
        CHECK los dominios simples que no cambian (franja horaria, tipo
        de prueba, tipo de novedad).
  \item \textbf{Resultado del intento.} El dominio de
        asignacion\_despacho.resultado queda en PENDIENTE, ENTREGADO,
        NOVEDAD y DEVUELTO, y el número de intento entre 1 y 3: los dos
        primeros son de entrega y el tercero es la devolución.
  \item \textbf{Integración futura.} Las columnas que dependen de otros
        módulos (envio.id\_factura, reclamacion.id\_factura) no se crean
        todavía; al integrar M2 se agregan sin modificar las columnas
        existentes.
\end{enumerate}

\subsubsection{Normalización y Formas Normales}

La normalización detecta y elimina redundancias que causan anomalías de
inserción, actualización y borrado. Una tabla está en primera forma
normal (1FN) si todos sus valores son atómicos; en segunda (2FN) si
está en 1FN y cada atributo no clave depende de la clave completa; en
tercera (3FN) si está en 2FN y ningún atributo no clave depende de otro
atributo no clave; y en forma normal de Boyce-Codd (FNBC) si todo
determinante es clave candidata.

Para comprobar el diseño se parte de cómo vería los datos alguien que
no ha modelado: una sola tabla con todo lo que se anota de una guía.

\begin{sloppypar}\noindent\textbf{Forma no normalizada.} GUIA(\pk{numero\_guia}, fecha\_registro,
nit\_cliente, razon\_social, direccion\_sede, zona\_sede,
documento\_destinatario, nombre\_destinatario, telefono\_destinatario,
direccion\_destino, zona\_destino, centro\_zona, modalidad,
plazo\_modalidad, peso\_declarado, peso\_real, contenido,
valor\_declarado, fragil, alto\_valor, valor\_flete, fecha\_compromiso,
\{fecha\_hora\_evento, estado, documento\_empleado, nombre\_empleado\},
\{id\_despacho, fecha\_despacho, placa, tipo\_vehiculo, capacidad\_tipo,
documento\_mensajero, nombre\_mensajero, orden\_visita, intento,
resultado, quien\_recibe, tipo\_prueba, tipo\_novedad,
descripcion\_novedad\}). Las llaves marcan grupos repetitivos: una guía
tiene varios eventos y varias visitas.\end{sloppypar}

\begin{sloppypar}\noindent\textbf{Primera forma normal.} Los grupos repetitivos salen a tablas
propias con la clave de la guía:
EVENTO(\pk{numero\_guia, fecha\_hora}, estado, documento\_empleado,
nombre\_empleado) y VISITA(\pk{numero\_guia, id\_despacho},
fecha\_despacho, placa, tipo\_vehiculo, capacidad\_tipo,
documento\_mensajero, nombre\_mensajero, orden\_visita, intento,
resultado, quien\_recibe, tipo\_prueba, tipo\_novedad,
descripcion\_novedad).\end{sloppypar}

\begin{sloppypar}\noindent\textbf{Segunda forma normal.} En VISITA, fecha\_despacho, placa,
tipo\_vehiculo, capacidad\_tipo y los datos del mensajero dependen solo
de id\_despacho, que es parte de la clave. Se separan en
DESPACHO(\pk{id\_despacho}, fecha, placa, tipo\_vehiculo,
capacidad\_tipo, documento\_mensajero, nombre\_mensajero). GUIA y
EVENTO ya cumplen 2FN: GUIA tiene clave simple y en EVENTO todos los
atributos dependen de la guía y del momento juntos.\end{sloppypar}

\textbf{Tercera forma normal.} Se eliminan las dependencias
transitivas:

\begin{itemize}[leftmargin=1.2cm,itemsep=1pt]
  \item En GUIA: documento\_destinatario $\rightarrow$ nombre, teléfono
        (sale DESTINATARIO); la guía pertenece a una solicitud, la
        solicitud a una sede y la sede a un cliente: sede $\rightarrow$
        cliente, dirección, zona y nit $\rightarrow$ razón social (salen
        SOLICITUD\_RECOLECCION, SEDE\_CLIENTE y CLIENTE); zona
        $\rightarrow$ centro (salen ZONA y CENTRO\_DISTRIBUCION);
        modalidad $\rightarrow$ plazo (sale MODALIDAD).
  \item En EVENTO: documento\_empleado $\rightarrow$ nombre (sale
        EMPLEADO) y estado $\rightarrow$ es\_final (sale ESTADO\_ENVIO).
  \item En DESPACHO: placa $\rightarrow$ tipo, marca, capacidad y tipo
        $\rightarrow$ capacidad máxima (salen VEHICULO y TIPO\_VEHICULO);
        documento\_mensajero $\rightarrow$ nombre (sale MENSAJERO, con
        sus datos de empleado en EMPLEADO).
  \item En VISITA no hay dependencias transitivas, pero los datos de
        prueba y de novedad son excluyentes y casi siempre nulos; por la
        decisión 3 de optimización salen a PRUEBA\_ENTREGA y NOVEDAD, y
        VISITA toma el nombre ASIGNACION\_DESPACHO con una clave
        sustituta.
\end{itemize}

\textbf{Forma normal de Boyce-Codd.} En las tablas resultantes todo
determinante es clave candidata. Las únicas tablas con más de una clave
candidata son DESTINATARIO, cuyo documento también la identifica, y
ASIGNACION\_DESPACHO, que también se identifica por (guía, despacho) y
por (guía, intento); ambas se declaran únicas. El resultado coincide con
las tablas obtenidas por conversión del modelo ER, lo que confirma que
el diseño conceptual era correcto.

% (Tabla de dependencias funcionales por tabla omitida en esta copia de
% trabajo: todas las tablas quedan en FNBC.)

\subsubsection{Diseño Lógico de la Base de Datos -- Modelo Relacional}
\label{sec:esquemarel}

La figura~\ref{fig:dmrelacional} presenta el modelo relacional del
módulo, construido en Data Modeler. Tiene 19 tablas con sus 19 claves
primarias, las 13 restricciones de unicidad y las 27 claves foráneas con
sus nombres.

En notación de esquema relacional (clave primaria subrayada, claves
foráneas en cursiva):

{\singlespacing\small\sloppy
\begin{itemize}[leftmargin=0.6cm,itemsep=3pt,label={}]
  \item \textbf{CENTRO\_DISTRIBUCION}(\pk{id\_centro}, nombre, direccion, telefono, estado)
  \item \textbf{ZONA}(\pk{id\_zona}, \fk{id\_centro}, nombre, localidades)
  \item \textbf{CLIENTE}(\pk{id\_cliente}, nit, razon\_social, tipo\_cliente, telefono, correo, fecha\_vinculacion, estado)
  \item \textbf{SEDE\_CLIENTE}(\pk{id\_sede}, \fk{id\_cliente}, \fk{id\_zona}, nombre, direccion, contacto, telefono)
  \item \textbf{MODALIDAD}(\pk{id\_modalidad}, nombre, plazo\_dias\_habiles)
  \item \textbf{TIPO\_VEHICULO}(\pk{id\_tipo\_vehiculo}, nombre, capacidad\_max\_kg)
  \item \textbf{VEHICULO}(\pk{id\_vehiculo}, \fk{id\_centro}, \fk{id\_tipo\_vehiculo}, placa, marca, modelo, capacidad\_kg, electrico, estado)
  \item \textbf{EMPLEADO}(\pk{id\_empleado}, \fk{id\_centro}, documento, nombres, apellidos, cargo, telefono, correo, fecha\_ingreso, fecha\_retiro, motivo\_retiro, estado)
  \item \textbf{MENSAJERO}(\pk{\fk{id\_empleado}}, \fk{id\_zona}, \fk{id\_tipo\_vehiculo})
  \item \textbf{DESTINATARIO}(\pk{id\_destinatario}, tipo\_documento, numero\_documento, nombre, telefono, correo)
  \item \textbf{DESPACHO}(\pk{id\_despacho}, \fk{id\_centro}, \fk{id\_vehiculo}, \fk{id\_mensajero}, fecha, tipo, hora\_salida, hora\_cierre, estado)
  \item \textbf{SOLICITUD\_RECOLECCION}(\pk{id\_solicitud}, \fk{id\_sede}, \fk{id\_despacho}, fecha\_solicitud, fecha\_programada, franja\_horaria, estado)
  \item \textbf{ENVIO}(\pk{numero\_guia}, \fk{id\_solicitud}, \fk{id\_destinatario}, \fk{id\_zona}, \fk{id\_modalidad}, fecha\_registro, direccion\_destino, peso\_declarado, peso\_real, contenido, valor\_declarado, fragil, alto\_valor, valor\_flete, fecha\_compromiso)
  \item \textbf{ASIGNACION\_DESPACHO}(\pk{id\_asignacion}, \fk{numero\_guia}, \fk{id\_despacho}, numero\_intento, orden\_visita, hora\_visita, resultado)
  \item \textbf{ESTADO\_ENVIO}(\pk{id\_estado}, nombre, es\_final)
  \item \textbf{EVENTO\_ENVIO}(\pk{id\_evento}, \fk{numero\_guia}, \fk{id\_estado}, \fk{id\_empleado}, fecha\_hora, observacion)
  \item \textbf{PRUEBA\_ENTREGA}(\pk{id\_prueba}, \fk{id\_asignacion}, fecha\_hora, nombre\_receptor, documento\_receptor, tipo\_prueba, archivo\_soporte)
  \item \textbf{NOVEDAD}(\pk{id\_novedad}, \fk{id\_asignacion}, tipo\_novedad, descripcion, accion, solucion, estado)
  \item \textbf{RECLAMACION}(\pk{id\_reclamacion}, \fk{id\_cliente}, \fk{numero\_guia}, fecha\_radicacion, tipo, descripcion, respuesta, fecha\_respuesta, estado)
\end{itemize}
\par}

\paragraph{Vistas de usuario (esquema externo)}

Cada usuario del módulo accede a los datos mediante vistas que reúnen lo
que necesita sin exponerle el resto (tabla~\ref{tab:vistas}). Las vistas
también implementan el estado actual del envío (decisión 2 de
optimización) y los indicadores de operación. Por ejemplo, la vista
del estado actual reúne cada envío con su cliente, su zona y su
modalidad, y toma como estado el de su evento más reciente.

{\singlespacing\small
\begin{longtable}{L{3.6cm} L{1.4cm} L{5.6cm} L{3.8cm}}
\caption{Vistas de usuario del módulo (esquema externo).}
\label{tab:vistas}\\
\toprule
\textbf{Vista} & \textbf{Usu.} & \textbf{Propósito} & \textbf{Tablas base} \\
\midrule
\endfirsthead
\toprule
\textbf{Vista} & \textbf{Usu.} & \textbf{Propósito} & \textbf{Tablas base} \\
\midrule
\endhead
\bottomrule
\endfoot
v\_seguimiento\_envio & U3, U9 & Recorrido completo de un envío por número de guía. El cliente (U9) la consulta filtrada por su id\_cliente. & envio, solicitud\_recoleccion, sede\_cliente, evento\_envio, estado\_envio, empleado \\
\midrule
v\_estado\_actual\_envio & U1, U2, U3 & Estado actual de cada envío (el de su último evento), con cliente, zona, modalidad y fecha de compromiso. & envio, solicitud\_recoleccion, sede\_cliente, cliente, zona, modalidad, evento\_envio, estado\_envio \\
\midrule
v\_planilla\_despacho & U2, U4 & Planilla de ruta: envíos de cada despacho en orden de visita, con datos del destinatario y del mensajero. & despacho, asignacion\_despacho, envio, destinatario, zona, vehiculo, empleado \\
\midrule
v\_envios\_por\_despachar & U2 & Envíos clasificados que todavía no están en un despacho de entrega abierto, por zona y modalidad, para armar las rutas. & v\_estado\_actual\_envio, envio, asignacion\_despacho, despacho \\
\midrule
v\_novedades\_abiertas & U3 & Novedades sin cerrar, con el intento en que ocurrieron y los datos para contactar al destinatario y al cliente. & novedad, asignacion\_despacho, envio, destinatario, solicitud\_recoleccion, sede\_cliente, cliente \\
\midrule
v\_reclamaciones\_pendientes & U1, U3 & Reclamaciones abiertas, con el cliente y el envío reclamado; insumo del indicador de tiempo de respuesta. & reclamacion, cliente \\
\midrule
v\_cumplimiento\_entregas & U1 & Por centro y zona: envíos entregados y cuántos se entregaron dentro de la fecha de compromiso (indicador de OE1). & envio, zona, evento\_envio, estado\_envio \\
\end{longtable}
\par}

\subsubsection{Documentación de la BD -- Diccionario de Datos}
\label{sec:diccionario}

% Copia de referencia: el diccionario completo (130
% columnas: tipo, nulo, clave, valor por defecto y descripción) es
% idéntico al DDL de Data Modeler en docs/datamodeler/, que trae cada
% columna con su tipo, NOT NULL, DEFAULT y COMMENT. Usar el DDL como
% fuente del diccionario.

Los estados del envío del catálogo ESTADO\_ENVIO son REGISTRADO,
RECOLECTADO, EN\_CENTRO, CLASIFICADO, EN\_RUTA, ENTREGADO, CON\_NOVEDAD,
DEVUELTO y ANULADO. Las modalidades son MISMO\_DIA, DIA\_SIGUIENTE y
ESTANDAR; los tipos de vehículo, MOTOCICLETA, FURGONETA y
BICICLETA\_ELECTRICA.
